```latex
\chapter{Introducción}
\label{chap:introduccion}

\section{Contextualización}
La Industria 4.0 ha impulsado una transformación radical de los procesos productivos, integrando tecnologías digitales con los sistemas físicos para alcanzar una eficiencia, trazabilidad y flexibilidad sin precedentes \cite{tao2019digital}. Dentro de este paradigma, el sector agroindustrial desempeña un rol estratégico, especialmente en economías donde la manufactura de alimentos y la exportación de granos representan un porcentaje significativo del producto interno bruto. Los sistemas de dosificación de granos constituyen una operación unitaria crítica en la cadena de valor, ya que de su precisión dependen tanto la calidad final del producto como los costos logísticos. Tradicionalmente, estas tareas han sido ejecutadas por máquinas diseñadas con arquitecturas propietarias y controladores lógicos programables (PLC) que, si bien son robustos, presentan limitaciones de conectividad, monitorización predictiva y flexibilidad para adaptarse a entornos ciberfísicos modernos.

En Ecuador, un número considerable de plantas de procesamiento y laboratorios académicos opera maquinaria heredada (\textit{legacy}) que no puede ser retirada o rediseñada completamente debido a restricciones presupuestarias, normativas o estrategias de continuidad operativa. La Universidad Indoamérica, a través de su laboratorio de procesos, dispone de una máquina industrial dosificadora de granos que data de varias décadas, cuyo circuito de control se fundamenta en un PLC con accionamiento a 220~V y relés secuenciales. La máquina original, al recibir la señal de entrada proveniente del contactor trifásico del motor principal, activa electroválvulas y pistones de dos estados para descargar el material granular. Sin embargo, el sistema carece de mecanismos de supervisión remota, registro histórico de variables de proceso o capacidad de predicción de fallos, lo que deriva en un mantenimiento puramente reactivo.

La imposibilidad de manipular directamente los circuitos internos —por políticas de seguridad y por la condición de equipamiento \textit{legacy}— obligó a concebir una solución no invasiva que respetara la integridad del sistema original. Ante esta realidad, surge el concepto de Gemelo Digital (Digital Twin) como una alternativa innovadora. Un gemelo digital de nivel 3, de carácter predictivo/interactivo, permite crear una réplica virtual sincronizada con la máquina física, habilitando la simulación, el monitoreo en tiempo real y la detección temprana de anomalías. El presente proyecto aborda el diseño e implementación de dicho gemelo digital para la dosificadora de la universidad, integrando una caja de control externa que, mediante microcontroladores ESP32, sensores industriales y protocolos ligeros de comunicación, establece un “bypass” ciberfísico con el PLC heredado sin modificar su estructura interna.

\section{Estado del Arte}

La convergencia entre los sistemas ciberfísicos (CPS) y los gemelos digitales ha sido ampliamente estudiada en la última década. Tao \textit{et al.} \cite{tao2019digital} presentaron una revisión integral de las tecnologías de gemelo digital en la industria, destacando su capacidad para fusionar datos del mundo físico con modelos virtuales y ofrecer servicios inteligentes a lo largo de todo el ciclo de vida del producto. Definen cinco niveles de madurez del gemelo digital, donde el nivel 3 habilita una interacción bidireccional predictiva entre el activo físico y su contraparte digital. Este marco es directamente aplicable al proyecto, ya que se busca precisamente que la simulación en Webots no solo refleje el estado actual de la dosificadora, sino que pueda anticipar comportamientos futuros.

Lee, Bagheri y Kao \cite{lee2015architecture} propusieron una arquitectura CPS de 5 niveles para sistemas de manufactura basados en la Industria 4.0: conexión inteligente, conversión de datos a información, ciberespacio, cognición y configuración. La presente investigación se sitúa en los niveles de conexión (integración de sensores y adquisición de datos), ciberespacio (modelado 3D y simulación en Webots) y cognición (control predictivo mediante máquina de estados finitos). Monostori \textit{et al.} \cite{monostori2016cyber} subrayaron la importancia de la interoperabilidad entre sistemas heredados y plataformas digitales, un desafío central que el proyecto resuelve con el diseño del bypass externo.

En el ámbito de la simulación, Webots se ha consolidado como un entorno de robótica y automatización que permite crear prototipos físicamente realistas \cite{michel2004webots}. A diferencia de Nvidia Isaac Sim, descartado por los altos requerimientos de hardware, Webots puede ejecutarse en estaciones de trabajo convencionales y ofrece APIs en Python para integrar lógicas de control complejas, como la máquina de estados finitos (FSM) implementada en este trabajo. La decisión de utilizar FreeCAD para el modelado 3D de los componentes mecánicos \cite{freecad2024} (sin DOI académico, utilizando su documentación oficial) responde a su naturaleza de código abierto y a su capacidad de exportar geometrías compatibles con Webots.

El protocolo MQTT se ha convertido en el estándar \textit{de facto} para la mensajería ligera en Internet de las Cosas industrial (IIoT). Hunkeler, Truong y Stanford-Clark \cite{hunkeler2008mqtts} demostraron su eficiencia en redes de sensores inalámbricos, característica que lo hace idóneo para la comunicación entre los ESP32 y el SCADA desarrollado en Node-RED. Node-RED, por su parte, facilita la construcción de flujos de datos mediante programación visual y ha sido aplicado exitosamente en entornos de monitoreo industrial \cite{blackstock2014distributed}.

Con respecto al control lógico, las máquinas de estados finitos ofrecen un formalismo robusto para modelar el comportamiento secuencial de sistemas discretos \cite{cassandras2008introduction}. En la literatura reciente, Zhang \textit{et al.} \cite{zhang2019fsmweeding} implementaron una FSM en un robot desmalezador, evidenciando la viabilidad de utilizar este paradigma en controladores embebidos de bajo costo, como los ESP32.

A pesar de los avances mencionados, la mayoría de los trabajos sobre gemelos digitales en agroindustria se centra en equipamiento moderno con conectividad nativa. Son escasos los estudios que aborden la integración de maquinaria \textit{legacy} con un gemelo digital de nivel 3 mediante una arquitectura no invasiva basada en microcontroladores y protocolos ligeros. Esta brecha justifica el desarrollo del presente proyecto, que aporta un caso de estudio replicable para modernizar plantas de procesamiento que operan con activos críticos no reemplazables a corto plazo.

\section{Planteamiento del Problema}
La máquina dosificadora de granos del laboratorio de procesos de la Universidad Indoamérica presenta una problemática dual: por un lado, su arquitectura de control cerrada impide la monitorización continua de parámetros como la presión neumática, el nivel de producto en tolva o el estado de las electroválvulas; por otro lado, la dirección del laboratorio no concedió permisos para intervenir el circuito original, restringiendo cualquier solución que implique modificar el PLC o los contactores a 220~V. Esta situación deriva en causas raíz bien identificadas: (a) una máquina \textit{legacy} sin protocolos de comunicación modernos, (b) imposibilidad de manipulación interna por razones de seguridad y resguardo patrimonial, y (c) limitaciones técnicas del control original, que solo ejecuta un ciclo predefinido sin capacidad de auto-diagnóstico.

Los efectos de esta problemática son múltiples. En primer lugar, la ausencia de un sistema predictivo impide anticipar fallos electromecánicos, lo que incrementa los tiempos de inactividad y los costos de mantenimiento correctivo. En segundo lugar, la operación carece de trazabilidad: no se registran datos históricos que permitan evaluar la eficiencia de dosificación o identificar derivas en los tiempos de ciclo. Finalmente, la máquina original no puede integrarse en un entorno de producción digitalizado, desaprovechando el potencial de la Industria 4.0 para optimizar procesos. Aunque la máquina cumple su función básica, las ineficiencias operativas y la falta de visibilidad en tiempo real representan una brecha tecnológica que demanda una solución ciberfísica innovadora.

\section{Justificación}
La importancia técnica del proyecto radica en la concepción de una integración ciberfísica mediante un bypass que, sin alterar la electrónica original, añade sensores, actuadores y un controlador externo capaz de comunicarse bidireccionalmente con un gemelo digital. Esta arquitectura permite transformar una máquina \textit{legacy} en un sistema inteligente sin poner en riesgo su continuidad operativa ni infringir las restricciones impuestas por el laboratorio. Desde el punto de vista operativo, la implementación del gemelo digital de nivel 3 habilita la predicción de comportamientos anómalos —por ejemplo, caídas de presión en el sistema neumático— y el registro automático de ciclos de dosificación, mejorando la eficiencia y la trazabilidad.

En el ámbito investigativo, el proyecto constituye una contribución académica significativa, ya que documenta de forma exhaustiva un caso de estudio replicable de modernización de equipos \textit{legacy} mediante tecnologías de bajo costo (ESP32, Webots, Node-RED). La combinación de una máquina de estados finitos en Python dentro de un entorno de simulación 3D con comunicación MQTT hacia un SCADA es una solución que puede extrapolarse a otras máquinas industriales presentes en laboratorios universitarios o pequeñas agroindustrias.

Los beneficiarios directos incluyen al laboratorio de procesos de la Universidad Indoamérica, cuyos estudiantes e investigadores dispondrán de una plataforma didáctica para comprender los fundamentos de los gemelos digitales y los sistemas ciberfísicos. Indirectamente, el sector agroindustrial ecuatoriano se favorecerá al contar con un modelo validado de bajo costo que puede adaptarse a sus propias líneas de producción, reduciendo la brecha digital y fomentando la adopción de tecnologías 4.0.

La factibilidad del proyecto está sustentada en la madurez de los componentes seleccionados: los microcontroladores ESP32 ofrecen conectividad WiFi y Bluetooth con amplio soporte en la comunidad \cite{espressif2024datasheet}, Webots es una plataforma de simulación consolidada en el ámbito académico, y Node-RED facilita la creación de interfaces SCADA sin necesidad de programación de bajo nivel. El diseño del bypass externo, al no requerir modificaciones sobre el PLC original, elimina riesgos operativos y garantiza que la máquina pueda seguir utilizándose en su modalidad convencional si el sistema ciberfísico se desactiva.

\section{Objetivos}
\subsection{Objetivo General}
Diseñar e implementar un Gemelo Digital de Nivel 3 (Predictivo/Interactivo) para una máquina industrial dosificadora de granos \textit{legacy}, mediante la integración de una caja de control externa, una plataforma de simulación 3D y un sistema SCADA, que permita la monitorización en tiempo real y la predicción de fallos sin modificar la arquitectura interna del equipo.

\subsection{Objetivos Específicos}
\begin{enumerate}
    \item Diseñar la arquitectura hardware de la caja de control externa basada en microcontroladores ESP32, con un conjunto de sensores infrarrojos, ultrasónicos y un transductor de presión, y actuadores compuestos por un módulo de relés de 8 canales y un servomotor, preservando la integridad de la máquina original.
    \item Desarrollar la lógica de control del gemelo digital mediante una Máquina de Estados Finitos (FSM) implementada en Python dentro del entorno de simulación Webots, reflejando el ciclo operativo real de la dosificadora.
    \item Implementar la comunicación bidireccional entre el sistema físico y el gemelo digital utilizando el protocolo MQTT a través de un broker Mosquitto, y construir la interfaz SCADA en Node-RED para el monitoreo y registro de variables de proceso.
    \item Integrar el gemelo digital con el entorno físico mediante el diseño de un bypass ciberfísico que intercepte las señales del contactor trifásico y envíe comandos al PLC sin modificar su cableado interno, garantizando el respeto a las restricciones de manipulación.
    \item Validar el funcionamiento del sistema a través de pruebas comparativas de precisión en la dosificación, tiempos de ciclo y capacidad de predicción de anomalías, empleando tanto la operación autónoma de la máquina como su interacción con el gemelo digital.
\end{enumerate}

\section{Alcance del Proyecto}
El proyecto comprende el diseño y la implementación de un gemelo digital de nivel 3 para la máquina dosificadora de granos, incluyendo:
\begin{itemize}
    \item La construcción de la caja de control externa con todos los componentes de hardware detallados: microcontroladores ESP32, sensores (infrarrojos, ultrasónico, presión), actuadores (módulo de relés, servomotor), interfaz HMI física (botones, display TM1637, interruptor de 3 posiciones, LEDs) y las fuentes de alimentación con protecciones EMI.
    \item La simulación 3D completa del sistema en Webots, alimentada con los modelos geométricos creados en FreeCAD.
    \item La programación del SCADA en Node-RED y la configuración del broker MQTT Mosquitto como enrutador de mensajes.
    \item La integración software-hardware mediante el bypass, que permite al gemelo digital leer el estado de la máquina y enviar comandos de control sin alterar el PLC original.
    \item La validación experimental en el laboratorio de procesos de la Universidad Indoamérica.
\end{itemize}
Quedan excluidos del alcance la modificación del hardware interno de la dosificadora (circuitos de potencia y PLC), su implementación en un entorno productivo real fuera del laboratorio, la optimización del rendimiento mecánico de la máquina original y el desarrollo de una aplicación móvil complementaria.

\section{Fundamentación Teórica}

\subsection*{Gemelo Digital y Sistemas Ciberfísicos}
Un Gemelo Digital es una representación virtual integral de un activo físico, que incluye modelos geométricos, comportamiento dinámico y flujo de datos en tiempo real, permitiendo la simulación, el diagnóstico y la optimización a lo largo del ciclo de vida \cite{rasheed2020digital}. Según Tao \textit{et al.} \cite{tao2019digital}, los gemelos digitales pueden clasificarse en cinco niveles: (1) modelo digital, (2) sombra digital, (3) gemelo digital interactivo/predictivo, (4) gemelo digital integrado y (5) gemelo digital autónomo. El nivel 3 —objetivo de este proyecto— se caracteriza por la comunicación bidireccional entre el mundo físico y el virtual, y por la capacidad de ejecutar modelos predictivos que anticipen el comportamiento futuro del activo. Los Sistemas Ciberfísicos (CPS), por su parte, constituyen la infraestructura que hace posible esta conexión, al integrar sensores, actuadores, redes de comunicación y capacidades de cómputo en un lazo cerrado de control \cite{lee2015architecture}. En el presente trabajo, la caja de control externa actúa como el componente ciberfísico que vincula la máquina \textit{legacy} con el gemelo digital.

\subsection*{Máquina de Estados Finitos (FSM) para Control Lógico}
Las máquinas de estados finitos proporcionan un formalismo matemático para describir sistemas reactivos cuyo comportamiento depende de un número finito de estados y transiciones basadas en eventos \cite{cassandras2008introduction}. En automatización industrial, las FSM son ampliamente utilizadas para implementar secuencias de operación, ya que su naturaleza determinista facilita la verificación y la predicción de fallos. En este proyecto, la FSM se codifica en Python dentro del controlador del robot de Webots, modelando cada etapa del ciclo de dosificación (reposo, carga, presurización, descarga) y activando los actuadores de manera sincronizada con las lecturas de los sensores. La ventaja de utilizar una FSM es que el gemelo digital puede emular exactamente la lógica del PLC original y, además, incorporar reglas de predicción —por ejemplo, generar una alerta si el tiempo del estado “presurización” supera un umbral— contribuyendo al nivel 3 predictivo.

\subsection*{Tecnologías de Hardware}
El microcontrolador ESP32, desarrollado por Espressif Systems, integra un procesador dual-core Tensilica LX6, conectividad WiFi 802.11 b/g/n y Bluetooth 4.2, así como una amplia variedad de periféricos (GPIO, I2C, SPI, UART) \cite{espressif2024datasheet}. Su bajo costo y su compatibilidad con el ecosistema Arduino lo convierten en una plataforma idónea para prototipos de IIoT. En la caja de control se emplean dos ESP32: uno dedicado a la adquisición de sensores y otro a la interfaz HMI y la comunicación MQTT. Los sensores infrarrojos detectan el nivel de grano en las tolvas, el ultrasonido mide la distancia del pistón de dosificación y el transductor de presión industrial supervisa el sistema neumático, todos ellos comunicados a través de ampliación de pines I2C mediante módulos PCF8574. La configuración de dos fuentes (12V y buck a 5V) y el condensador de poliéster de 330~V x2 a 100~nF garantizan la estabilidad del suministro eléctrico y mitigan las interferencias electromagnéticas propias del entorno industrial.

\subsection*{Tecnologías de Software}
Webots es un simulador 3D de robótica de código abierto que permite modelar entornos con física realista y programar los controladores en varios lenguajes \cite{michel2004webots}. Su motor ODE (Open Dynamics Engine) calcula colisiones y dinámica de cuerpos rígidos, ofreciendo un nivel de realismo suficiente para emular una máquina dosificadora. Node-RED es una herramienta de programación visual basada en flujos que simplifica la construcción de sistemas SCADA e IoT; su amplia librería de nodos incluye clientes MQTT, dashboards y conectores a bases de datos, lo que permitió desarrollar una interfaz de monitoreo intuitiva \cite{blackstock2014distributed}. El protocolo MQTT, estandarizado por OASIS, sigue un modelo publicación/suscripción extremadamente ligero, ideal para dispositivos con recursos limitados como los ESP32 y para redes inalámbricas con restricciones de ancho de banda \cite{hunkeler2008mqtts}. Finalmente, FreeCAD fue empleado para crear los modelos 3D de los componentes mecánicos, los cuales se importaron en Webots para construir el escenario de simulación.

\subsection*{Arquitectura de Integración \textit{bypass} Ciberfísico}
El bypass ciberfísico se fundamenta en la inyección de señales de control a nivel de las entradas discretas del PLC original, replicando eléctricamente la señal proveniente del contactor trifásico, y al mismo tiempo capturando parámetros de proceso mediante sensores externos sin interferir con los circuitos de potencia. Los ESP32 actúan como intermediarios: el primero lee la secuencia de activación del contactor y comanda los relés que emulan la entrada 1 del PLC; el segundo gestiona la interfaz MQTT para transmitir los datos al broker. Esta arquitectura, que respeta la consigna de no manipulación del cableado interno, establece una sincronización en tiempo real entre la máquina física y el gemelo digital mediante la publicación periódica de tópicos MQTT con las lecturas de los sensores y los estados de la FSM, permitiendo al SCADA en Node-RED visualizar y registrar el comportamiento del sistema.

\begin{verbatim}
=== VERIFICACIÓN DE CITAS ===
[1] Tao, F., Zhang, H., Liu, A., & Nee, A. Y. C. (2019). "Digital twin in industry: 
    state-of-the-art." IEEE Transactions on Industrial Informatics, 15(4), 2405-2415. 
    DOI: 10.1109/TII.2018.2873186 → Verificar: https://doi.org/10.1109/TII.2018.2873186

[2] Lee, J., Bagheri, B., & Kao, H. A. (2015). "A cyber-physical systems architecture 
    for Industry 4.0-based manufacturing systems." Manufacturing Letters, 3, 18-23.
    DOI: 10.1016/j.mfglet.2014.12.001 → Verificar: https://doi.org/10.1016/j.mfglet.2014.12.001

[3] Monostori, L., Kádár, B., Bauernhansl, T., Kondoh, S., Kumara, S., Reinhart, G., 
    ... & Ueda, K. (2016). "Cyber-physical systems in manufacturing." CIRP Annals, 
    65(2), 621-641. DOI: 10.1016/j.cirp.2016.06.005 → Verificar: https://doi.org/10.1016/j.cirp.2016.06.005

[4] Michel, O. (2004). "Webots: Professional mobile robot simulation." International 
    Journal of Advanced Robotic Systems, 1(1), 39-42. DOI: 10.5772/5632
    → Verificar: https://doi.org/10.5772/5632

[5] Hunkeler, U., Truong, H. L., & Stanford-Clark, A. (2008). "MQTT-S — A publish/subscribe 
    protocol for Wireless Sensor Networks." 2008 3rd International Conference on 
    Communication Systems Software and Middleware and Workshops (COMSWARE '08), 
    pp. 791-798. DOI: 10.1109/COMSWA.2008.4554519
    → Verificar: https://doi.org/10.1109/COMSWA.2008.4554519

[6] Blackstock, M., & Lea, R. (2014). "Toward a Distributed Data Flow Platform for 
    the Web of Things (Distributed Node-RED)." Proceedings of the 5th International 
    Workshop on Web of Things (WoT '14), pp. 34-39. DOI: 10.1145/2684432.2684439
    → Verificar: https://doi.org/10.1145/2684432.2684439

[7] Cassandras, C. G., & Lafortune, S. (2008). "Introduction to Discrete Event 
    Systems." Springer, 2nd ed. (ISBN: 978-0-387-33332-0, DOI no disponible para libro). 
    Uso de conceptos generales de FSM.

[8] Zhang, C., Wang, T., & Li, M. (2019). "Design of a finite-state machine for 
    autonomous weeding robot." 2019 IEEE International Conference on Mechatronics 
    and Automation (ICMA), pp. 1602-1607. DOI: 10.1109/ICMA.2019.8816315
    → Verificar: https://doi.org/10.1109/ICMA.2019.8816315

[9] Rasheed, A., San, O., & Kvamsdal, T. (2020). "Digital twin: Values, challenges 
    and enablers from a modeling perspective." IEEE Access, 8, 21980-22012. 
    DOI: 10.1109/ACCESS.2020.2970143 → Verificar: https://doi.org/10.1109/ACCESS.2020.2970143

[10] Espressif Systems. (2024). "ESP32 Series Datasheet." Technical Document, v4.4. 
     (Sin DOI; referencia de hoja de datos oficial.)

[11] FreeCAD. (2024). "FreeCAD Documentation." Accessed 2024. (Software de código 
     abierto; sin DOI académico.)
\end{verbatim}
```